\section{Introduction}

Large language models are trained on massive human language corpora,
which induce a characteristic geometry in learned representations.
Across layers and training stages, activations, gradients, and model parameters
exhibit pronounced anisotropy:
a small number of directions concentrate disproportionate energy
(strong spectral spikes),
while the remaining directions form a broad long tail
\cite{ethayarajh2019contextual,mu2018all,timkey2021rogue,kovaleva2021bertbusters,machina2024anisotropy}.

This anisotropy is not accidental.
It reflects the hierarchical structure of human language itself.
At every level of representation—tokens, phrases, clauses, and full sentences—
there exist shared semantic components that are common across tokens within a context.
When contextual aggregation mechanisms such as attention and nonlinear transformations
extract these shared semantics, they naturally induce coherent mean directions
in the representation space.
Consequently, mean bias is not confined to the embedding layer;
it re-emerges across layers as models progressively capture higher-level
semantic regularities.

While such structure is linguistically natural,
it becomes a computational liability during large-scale training.
Anisotropy inflates elementwise dynamic range,
distorts adaptive second-moment estimates
\cite{kingma2014adam,loshchilov2017adamw,spectra2026},
and compresses long-tail variation in low-bit regimes
\cite{xiao2023smoothquant,frantar2022gptq,lin2024awq,li2024svdquant,ashkboos2024quarot}.
In high-dimensional spaces, even moderate coherent bias produces
a consistent directional drift across thousands of coordinates,
amplifying outliers and destabilizing numerical precision.
Thus anisotropy constitutes a structural \emph{curse}
for efficient and stable large-scale training.

In this work, we perform a systematic layer-wise investigation
of anisotropy across activations, gradient outputs, and model parameters.
We uncover a structural asymmetry:
\emph{mean bias is a dominant component of activation anisotropy across layers,
whereas gradient outputs exhibit no comparable coherent mean}.
Moreover, although model parameters display spectral spikes,
they do not accumulate persistent mean offsets.
This establishes activation mean bias as the principal forward-path driver
of anisotropic geometry.

Importantly, we show that anisotropy can be decomposed into two components:
(i) a coherent mean bias, and
(ii) a residual low-rank spike.
The mean bias is the primary source of extreme elementwise outliers
and adaptive second-moment inflation,
while the residual spike plays a secondary role.
Thus the dominant instability mechanism is not generic high-dimensional variance,
but a structured rank-one component that propagates through residual connections
across depth.

Tracing this structure backward,
we identify its statistical origin in the skewed token distribution of natural language.
Embedding weights develop frequency-aligned structure
\cite{mikolov2013efficient,mikolov2013distributed,arora2017simple},
and subsequent contextual aggregation amplifies shared semantics at progressively
higher abstraction levels.
Mean bias therefore reflects the hierarchical organization of world knowledge
encoded in language rather than an architectural accident.

Crucially, this structure is not merely harmful—it is also exploitable.
Because the dominant instability arises from a rank-one component,
it can be isolated and removed with a single mean-subtraction operation.
The remaining low-rank spike can, if necessary,
be whitened once at the embedding activation level.
Both operations are hardware-friendly and substantially cheaper
than repeated layer-wise SVD corrections.
In this sense, the same structure that creates instability
also enables efficient mitigation—the \emph{blessing} of bias.

We validate this bias-centric interpretation in two downstream settings.
In low-bit quantization, mean removal significantly reduces dynamic-range inflation
\cite{xiao2023smoothquant,wei2023outlier,ashkboos2024quarot}.
In adaptive optimization, it contracts extreme second-moment tails,
improving effective learning-rate allocation
\cite{kingma2014adam,spectra2026}.
These results suggest that addressing bias at its source
can replace or complement spectral interventions
with minimal computational overhead.

\paragraph{Contributions.}
\begin{enumerate}[leftmargin=*,itemsep=2pt]
    \item We show that activation mean bias emerges across layers as a structural consequence of hierarchical semantic aggregation.
    \item We demonstrate that mean-dominated components account for most extreme activation magnitudes and second-moment inflation.
    \item We establish a forward-path asymmetry: gradient outputs remain mean-neutral while activations accumulate coherent bias.
    \item We propose source-level mean removal and optional spike whitening as efficient spectral conditioning mechanisms.
    \item We demonstrate improvements in low-bit quantization and adaptive optimization.
\end{enumerate}